\par\Needspace{6\baselineskip}\medskip
\hypertarget{native-159-article114}{}\pdfbookmark[2]{114.}{article-114}
\begin{center}114.\end{center}
Multipliziert man Gleichung 1 mit $\sin\alpha''\operatorname{tang}B''-\sin L''\operatorname{tang}\beta''$, Gleichung 2 mit $\cos L''\operatorname{tang}\beta''-\cos\alpha''\operatorname{tang}B''$ und Gleichung 3 mit $\sin(L''-\alpha'')$ und addiert die Produkte, so erhält man
\[
\begin{aligned}
0={}&n\{(0.2.\mathrm{II})\delta+(\mathrm{O}.2.\mathrm{II})D\}\\
&-n'\{(1.2.\mathrm{II})\delta'+(\mathrm{I}.2.\mathrm{II})D'\}
\end{aligned}\tag*{4)}
\]
und ebenso oder, bequemer, durch bloßes Vertauschen der Orte
\[
\begin{aligned}
0={}&n\{(0.1.\mathrm{I})\delta+(\mathrm{O}.1.\mathrm{I})D\}\\
&+n''\{(2.1.\mathrm{I})\delta''+(\mathrm{II}.1.\mathrm{I})D''\}
\end{aligned}\tag*{5)}
\]
\[
\begin{aligned}
0={}&n'\{(1.0.\mathrm{O})\delta'+(\mathrm{I}.0.\mathrm{O})D'\}\\
&-n''\{(2.0.\mathrm{O})\delta''+(\mathrm{II}.0.\mathrm{O})D''\}.
\end{aligned}\tag*{6)}
\]
Ist also das Verhältnis der Größen $n$, $n'$ gegeben, so lässt sich mithilfe von Gleichung 4 aus $\delta$ der Wert von $\delta'$ oder der Wert von $\delta$ aus $\delta'$ bestimmen. Entsprechendes gilt für die Gleichungen 5 und 6. Aus der Verbindung der Gleichungen 4, 5 und 6 ergibt sich die Gleichung
\[
\begin{aligned}
&\frac{(0.2.\mathrm{II})\delta+(\mathrm{O}.2.\mathrm{II})D}{(0.1.\mathrm{I})\delta+(\mathrm{O}.1.\mathrm{I})D}\\[.6em]
&\qquad\times\frac{(1.0.\mathrm{O})\delta'+(\mathrm{I}.0.\mathrm{O})D'}{(1.2.\mathrm{II})\delta'+(\mathrm{I}.2.\mathrm{II})D'}\\[.6em]
&\qquad\times\frac{(2.1.\mathrm{I})\delta''+(\mathrm{II}.1.\mathrm{I})D''}{(2.0.\mathrm{O})\delta''+(\mathrm{II}.0.\mathrm{O})D''}=-1,
\end{aligned}\tag*{7)}
\]
mit der man aus zwei Abständen des Himmelskörpers von der Erde den dritten bestimmen kann. Es lässt sich jedoch zeigen, dass Gleichung 7 zu einer Identität wird und somit zur Bestimmung eines Abstands aus den beiden anderen ungeeignet ist, wenn $B=B'=B''=0$ und
\[
\left.\begin{aligned}
&\operatorname{tang}\beta'\operatorname{tang}\beta''\sin(L-\alpha)\sin(L''-L')\\
&+\operatorname{tang}\beta''\operatorname{tang}\beta\sin(L'-\alpha')\sin(L-L'')\\
&+\operatorname{tang}\beta\operatorname{tang}\beta'\sin(L''-\alpha'')\sin(L'-L)
\end{aligned}\right\}=0.
\]
\par
Die folgende Formel vermeidet diese Schwierigkeit und ergibt sich \SourceLeaf{160}leicht aus den Gleichungen 1, 2 und 3.
\[
\begin{aligned}
&(0.1.2)\delta\delta'\delta''+(\mathrm{O}.1.2)D\delta'\delta''\\
&+(0.\mathrm{I}.2)D'\delta\delta''+(0.1.\mathrm{II})D''\delta\delta'\\
&+(0.\mathrm{I}.\mathrm{II})D'D''\delta+(\mathrm{O}.1.\mathrm{II})DD''\delta'\\
&+(\mathrm{O}.\mathrm{I}.2)DD'\delta''+(\mathrm{O}.\mathrm{I}.\mathrm{II})DD'D''=0.
\end{aligned}\tag*{8)}
\]
\par
Multipliziert man Gleichung 1 mit $\sin\alpha'\operatorname{tang}\beta''-\sin\alpha''\operatorname{tang}\beta'$, Gleichung 2 mit $\cos\alpha''\operatorname{tang}\beta'-\cos\alpha'\operatorname{tang}\beta''$ und Gleichung 3 mit $\sin(\alpha''-\alpha')$ und addiert die Produkte, so erhält man
\[
\begin{aligned}
0={}&n\{(0.1.2)\delta+(\mathrm{O}.1.2)D\}\\
&-n'(\mathrm{I}.1.2)D'+n''(\mathrm{II}.1.2)D''
\end{aligned}\tag*{9)}
\]
und ebenso
\[
\begin{aligned}
0={}&n(0.\mathrm{O}.2)D-n'\{(0.1.2)\delta'+(0.\mathrm{I}.2)D'\}\\
&+n''(0.\mathrm{II}.2)D''
\end{aligned}\tag*{10)}
\]
\[
\begin{aligned}
0={}&n(0.1.\mathrm{O})D-n'(0.1.\mathrm{I})D'\\
&+n''\{(0.1.2)\delta''+(0.1.\mathrm{II})D''\}.
\end{aligned}\tag*{11)}
\]
Mithilfe dieser Gleichungen lassen sich aus dem bekannten Verhältnis der Größen $n$, $n'$, $n''$ die Abstände $\delta$, $\delta'$, $\delta''$ ermitteln. Dieser Schluss gilt jedoch nur im Allgemeinen und lässt eine Ausnahme zu, wenn $(0.1.2)=0$ ist. In diesem Fall folgt nämlich, wie sich zeigen lässt, aus den Gleichungen 9, 10 und 11 lediglich eine notwendige Beziehung zwischen den Größen $n$, $n'$, $n''$, und zwar aus jeder der drei Gleichungen dieselbe. Die entsprechenden Einschränkungen für die Gleichungen 4, 5 und 6 wird der kundige Leser ohne Weiteres erkennen.
\par
Im Übrigen sind alle hier entwickelten Schlussfolgerungen unbrauchbar, wenn die Bahnebene mit der Ekliptik zusammenfällt. Sind nämlich $\beta$, $\beta'$, $\beta''$, $B$, $B'$, $B''$ sämtlich $=0$, so ist Gleichung 3 eine \emph{Identität}, und damit sind es auch alle darauf folgenden Gleichungen.
\par\bigskip\begin{center}\rule{.23\linewidth}{.35pt}\end{center}
